\chapter*{Preface}
\addcontentsline{toc}{chapter}{Preface}

Digital systems are made of code, data, interfaces, organisations, and decisions. A program can be syntactically correct and still solve the wrong problem. A database can be fast and still encode a harmful policy. A beautiful interface can be inaccessible. A predictive model can be accurate on a test set and still be unsafe in use. This book is organised around those connections.

The intended reader may arrive from computer science, social science, design, teaching, business, or another discipline. The first chapters therefore begin with ordinary examples and gradually introduce formal language. The aim is not to lower the intellectual level. It is to make the reasoning visible before asking the reader to manipulate its notation.

The scope of this book was inspired by the broad range of questions encountered in a modern information-technology degree, but the book is not a degree programme, course catalogue, or institutional manual. It is a standalone synthesis. Its purpose is to show how ideas from computing, software engineering, interaction design, organisations, research, and artificial intelligence fit together when we try to understand and build digital systems.

One recurring case, CivicQueue, follows a fictional public-service appointment system from an uncertain problem to a tested prototype, a relational database, a deployable web service, an organisational intervention, and a bounded research evaluation. The case is deliberately imperfect. It gives readers opportunities to identify assumptions, privacy risks, competing stakeholder interests, and claims that the evidence cannot support.

The central question is simple: \emph{what must be true for this digital system to be a responsible answer to this problem?} The chapters provide different ways of answering it: computational models, algorithms, software architecture, interaction design, organisational analysis, empirical research, and critical reflection. You can read the book from beginning to end, or enter through the chapter that matches the problem in front of you.

\vfill
\begin{flushright}
July 2026
\end{flushright}
